\documentclass[line, margin, 10pt]{res}

\usepackage[T1]{fontenc}
\usepackage{enumitem}
\usepackage{hyperref}
\hypersetup{colorlinks=true, urlcolor=black, linkcolor=black}

% Widen the section label column slightly for multi-word labels
\newsectionwidth{1.1in}

% Tighten up list spacing
\setlist[itemize]{leftmargin=1.2em, itemsep=1pt, topsep=2pt, parsep=0pt}

% Page margins — res.cls sets textheight=9in by default; leave it alone
% Just nudge the side margin and cap textheight sensibly
\addtolength{\oddsidemargin}{-.1in}
\setlength{\textwidth}{5.6in}
\setlength{\textheight}{8.5in}
\addtolength{\topmargin}{.1in}

\begin{document}

% ---------------------------------------------------------
% Header: name left, rule, contact right
% ---------------------------------------------------------
\name{\Large Alexander Joseph Almeida}

\begin{resume}

% ---------------------------------------------------------
% Contact block as first section, two-column, matching Paciorek style
% ---------------------------------------------------------
\section{\sc Contact\\ Information}
\begin{ncolumn}{2}
  Graduate School of Business          & \textit{WWW:} \href{https://alexjalmeida.github.io}{alexjalmeida.github.io} \\
  Stanford University                  & \textit{E-mail:} \href{mailto:alexander.almeida@stanford.edu}{alexander.almeida@stanford.edu} \\
  655 Knight Way                       & \\
  Stanford, CA 94305                   & \\
\end{ncolumn}

% ---------------------------------------------------------
\section{\sc Fields of\\ Interest}
Econometrics, Causal Inference, Machine Learning

% ---------------------------------------------------------
\section{\sc Education}
% ---------------------------------------------------------

\textbf{Stanford University}, Stanford, CA \\
Ph.D., Economic Policy, Graduate School of Business \hfill 2024 -- 2030 (expected) \\
\hspace*{1em}\textit{Advisors:} Guido Imbens, Lihua Lei \\
M.S., Statistics \hfill 2024 -- 2027 (expected) \\
Research Fellow, Graduate School of Business \hfill 2022 -- 2024 \\[2pt]
\textit{Fellowships \& Awards:} EDGE Doctoral Fellowship (Stanford VPGE)

\vspace{4pt}

\textbf{Arizona State University} -- Barrett, The Honors College, Tempe, AZ \\
B.S.\ Economics \& B.S.\ Mathematics, GPA: 4.00/4.00 \hfill 2018 -- 2022 \\[2pt]
\textit{Honors:} Dean's Medal -- Economics (Most Outstanding Graduate),
Jonathan D.\ and Helen Wexler Award (Top 12 Seniors in Mathematics),
Moeur Award (8 semesters above 4.0), Phi Beta Kappa, Beta Gamma Sigma \\[2pt]
\textit{Scholarships:} National Hispanic Scholar (Full Tuition),
IGNITE Fellowship -- Teach for America, Grady Gammage Scholarship,
Earl \& Ellen Davis Memorial Scholarship, SCETL Founders' Scholarship,
Business School Scholarship

% ---------------------------------------------------------
\section{\sc Working\\ Papers}
% ---------------------------------------------------------

Almeida, Alexander, Susan Athey, Guido Imbens, Eva Lestant, and Alexia Olaizola.
``Estimating Variances for Causal Panel Data Estimators.''
\textit{arXiv:2510.11841}, October 2025.
\begin{itemize}
  \item Develops a unified framework for variance estimation in causal panel data
        settings; shows that the three main approaches (regression-based,
        unit-placebo, and time-placebo) each target a different conditional variance
        under distinct exchangeability assumptions, and proposes a new
        heteroskedasticity-robust estimator with superior statistical power
\end{itemize}

\vspace{2pt}

Almeida, Alexander, and Lihua Lei.
``Compound Selection with Heteroskedastic Errors of Unknown Variance.''
\textit{Work in progress.}
\begin{itemize}
  \item Develops an empirical Bayes approach to compound selection decisions,
        extending prior work by Lei et al.\ to settings where observation variances
        are unknown; relaxes the known-variance assumption in Stein's unbiased risk
        estimation to handle heteroskedastic data
\end{itemize}

% ---------------------------------------------------------
\section{\sc Research\\ Experience}
% ---------------------------------------------------------

\textbf{Stanford Graduate School of Business}, Stanford, CA \\
\textit{Research Fellow} -- supervised by Prof.\ Rebecca Diamond and Prof.\ Guido Imbens
\hfill Jun.\ 2022 -- Jun.\ 2024
\begin{itemize}
  \item Developed heteroskedasticity-robust variance estimators for causal panel
        data settings, unifying regression-based, unit-placebo, and time-placebo
        approaches under a common exchangeability framework
  \item Constructed and managed terabyte-scale administrative data pipelines
        supporting empirical research in labor and public economics
  \item Contributed to three publications in labor and public economics (Garrett, Ohrn, Su\'{a}rez Serrato; Curtis et al.; Su\'{a}rez Serrato, Wingender)
\end{itemize}

\vspace{4pt}

\textbf{Federal Reserve Bank of Dallas}, Dallas, TX \\
\textit{Research Intern} -- supervised by Dr.\ Ali Ozdagli
\hfill Jun.\ 2021 -- Aug.\ 2021
\begin{itemize}
  \item Authored two original blog posts for the Dallas Fed Economics blog drawing
        on research on firm financing and monetary policy transmission
  \item Developed a literature review on environmental risk in financial markets
\end{itemize}

\vspace{4pt}

\textbf{Future H2O Economic Research}, Tempe, AZ \\
\textit{Research Assistant} -- supervised by Prof.\ W.\ Michael Hanemann
\hfill Mar.\ 2021 -- May 2022
\begin{itemize}
  \item Investigated hedonic property valuation of ecological attributes in Texas
        flood basins
\end{itemize}

\vspace{4pt}

\textbf{College of Health Solutions: Behavioral Health Research Lab}, Tempe, AZ \\
\textit{Research Assistant} -- supervised by Prof.\ Chad Stecher
\hfill Sep.\ 2020 -- May 2022
\begin{itemize}
  \item Designed and implemented a randomized controlled experiment to evaluate
        behavioral automaticity
  \item Applied machine learning methods to large mobile-health databases
\end{itemize}

% ---------------------------------------------------------
\section{\sc Presentations}
\begin{tabular}{@{}p{0.45in}p{4.95in}@{}}
2026 & Empirical Methods in the Age of AI, Stanford University (poster)
\end{tabular}

% ---------------------------------------------------------
\section{\sc Teaching}
% ---------------------------------------------------------

\textbf{Arizona State University -- Economics Department}, Tempe, AZ \\
\textit{Teaching Assistant} -- Foundations of Microeconomics
\hfill Oct.\ 2019 -- Jan.\ 2021
\begin{itemize}
  \item Assisted in developing coursework for a distance-learning MOOC platform
        serving over 800 students
\end{itemize}

\vspace{4pt}

\textbf{Arizona State University Academic Success Programs}, Tempe, AZ \\
\textit{Tutor} -- Economics, Finance, Mathematics
\hfill Feb.\ 2020 -- Feb.\ 2021

% ---------------------------------------------------------
\section{\sc Skills}
% ---------------------------------------------------------

\textbf{Computing:} Bash, Python, R, Stata, Julia, MATLAB, C++, Java,
\LaTeX, SQL; some experience in PyTorch, Tableau \\[3pt]
\textbf{Languages:} English (native), Spanish (fluent)

% ---------------------------------------------------------
\section{\sc Service \&\\ Affiliations}
% ---------------------------------------------------------

\begin{itemize}
  \item \textbf{Gammage Scholarship Program} -- Chair and Project Leader,
        service-based leadership program for National Merit Scholars
        (2018--2022)
  \item \textbf{RISE Tutoring} -- Finance Committee Head, Refugee Integration
        Stabilization and Education Scholars; recipient of ASU Pitchfork Award
        for outstanding student organizations (2019--2022)
\end{itemize}

\end{resume}
\end{document}
